\section{Conclusion}
\label{sec:conclusion}

LEO constellations forward IP traffic in orbit across many on-board routers
whose compute and memory cannot sustain conventional per-flow security, yet whose
open ISLs demand protection. 6FLAS resolves this tension by
moving security state out of the satellite and into the packet: it derives a
domain-authenticated flow identifier from a domain key and carries it in the immutable
IPv6 Flow Label, binds it to an SRv6 Security SID that encodes the per-hop path
policy, and verifies each packet with a single keyed MAC over the authenticated routing context plus a constant-time,
topology-aware egress check. The design needs no per-flow security association,
so it survives ISL handovers without re-keying and resists state-exhaustion
denial of service by construction.

6FLAS provides domain-origin authentication (not per-sender accountability within a domain), routing-context integrity,
probabilistic anti-replay within the window lifetime (with long-term mitigation via key rotation), and access control. The forgery probability is bounded by the MAC security and the $32$-bit tag truncation. The evaluation confirms that 6FLAS cuts per-hop processing by
$6.0\times$ versus an SRv6 HMAC baseline operating over a similar authenticated span, and by $55\times$ versus IPsec ESP (which additionally encrypts payload); empirical measurements
confirm the structural ordering $\text{FLAS} \ll \text{SRv6-Sec} \ll \text{IPsec}$.
Header overhead drops by $76\%$,
per-flow state is eliminated, and end-to-end latency remains
propagation-dominated. Future work will
add a composable confidentiality layer and
a standardization path for the Flow-Label security profile and the
\texttt{End.SEC} behavior.

As mega-constellations scale to tens of thousands of satellites, stateless security mechanisms become essential for sustainable on-board operations.
